% !TEX root = ../main.tex
\chapter{Option Prices and Probability Duality}\label{ch:duality}

\begin{sources}
\small
\begin{tabularx}{\linewidth}{@{}l l L@{}}
\toprule
\textbf{Lecture} & \textbf{Speaker} & \textbf{Content}\\
\midrule
2013 L20 & S.~Blythe (Harvard Management Company) & Guest lecture, transcript only (no slides, no notes on OCW). Three payoffs (call, zero-coupon bond, digital); the call spread as a discrete first strike-derivative and its limit, the digital; FTAP and the risk-neutral density as the second strike-derivative of the call price; the call butterfly as a traded proxy for the density; the ``other direction'' (density $\to$ price via FTAP/LOTUS); calls as a spanning set for European payoffs, first by piecewise-linear replication and then by an exact second-order Taylor identity; a remark that the terminal distribution for every maturity does not pin down the stock's stochastic process (Dupire--Derman), with the practical consequence for exotic-option pricing.\\
\bottomrule
\end{tabularx}
\smallskip
\textbf{Files.} None: OCW lists no slides or lecture notes for L20, only the video transcript.
\begin{coursenote}[2013 L22 has no material at all]
The following lecture, 2013 L22, ``Calculus of Variations and its Application in FX Execution''
(E.~Pan, Morgan Stanley), has neither a transcript nor slides on OCW and is not covered anywhere in
these notes.
\end{coursenote}
\textbf{Problem sets.} None. The exercises at the end of the chapter are all ``(not from the course)''.
\textbf{Reconstruction.} Blythe states the Breeden--Litzenberger relation and the spanning formula but
proves neither in full generality (the spanning proof is done only for piecewise-linear $g$, by
inspection). \cref{ch22:sec-bl-recap} is a two-line recap of the proof already given in
\cref{ch03:cor-bl}. \cref{ch22:sec-butterfly} adds the triangle/rectangle area argument only sketched
in class. \cref{ch22:sec-spanning} gives the exact Taylor-with-remainder computation in full (the
lecture writes it down and calls it ``exact'', without proof, and skips the discretisation
$g\to$ piecewise-linear step). All numerics (\cref{ch22:tab-butterfly,ch22:fig-density,ch22:ex-x2})
are computed for these notes, with the Black--Scholes parameters of \cref{ch15:ex-atm}
($S_0=K=80$, $T=2$, $r=5\%$, $\sigma=20\%$) so that the digital and call-spread numbers reproduce
\cref{ch15:eq-digital-limit}--\eqref{ch15:eq-digital}.
\end{sources}

\section*{Motivation}
Chapter~\ref{ch:bs} already used one instance of a strike derivative: the digital option is
$-\partial C/\partial K$ (\cref{ch15:sec-digital}), and \cref{ch03:cor-bl} proved that a call price
viewed as a function of strike encodes the whole distribution of the underlying, $C''(K)=f(K)$. This
chapter is about what that fact \emph{means}. Blythe's title for it, ``option-probability duality'',
is the claim that a call option is not a separate, exotic object requiring a pricing model: it
\emph{is} a piece of probability information, tradeable and observable, and the map between ``the set
of all call prices'' and ``the risk-neutral distribution of the underlying'' goes both ways. Read one
way, market call prices reveal the market's implied probabilities (this is genuinely done by trading
desks, via the \emph{call butterfly} of \cref{ch22:sec-butterfly}). Read the other way, the Fundamental
Theorem of Asset Pricing (\cref{thm:ftap}, generalised here to a continuum of states in
\cref{ch22:sec-ftap}) turns any candidate distribution into a full price system for every European
payoff. The two directions meet in the \emph{spanning} result of \cref{ch22:sec-spanning}: knowing the
call price at every strike is exactly as much information as knowing the whole payoff-to-price map,
because every payoff is a specific, explicit portfolio of zero-coupon bond, stock and calls. The
chapter closes in \cref{ch22:sec-process} with what this does \emph{not} give: the terminal
distribution at every maturity is not enough to pin down the process the stock follows in between,
which is why option markets alone cannot price every exotic derivative and a model (\cref{ch:sde}) is
still needed for path-dependent payoffs.

% ==================================================================
\section{Three payoffs and one assumption}\label{ch22:sec-ftap}

Fix a maturity $T$ and let $S_T$ be the price of the underlying at $T$. Following the lecture
\ts{13}{20}{13:00}, consider three contracts, all defined purely by their payout at $T$:
\begin{itemize}
  \item a \emph{call} with strike $K$: payout $(S_T-K)^+$;
  \item a \emph{zero-coupon bond}: payout $1$;
  \item a \emph{digital} with strike $K$: payout $\ind\{S_T>K\}$.
\end{itemize}
Write $C_K(t,T)$, $Z(t,T)$, $D_K(t,T)$ for their prices at $t\le T$; $Z(T,T)=1$. As in
\cref{def:onepmodel}, a payoff that is a deterministic function of $S_T$ alone is called a
\emph{European} claim, written $g(S_T)$ for a payout function $g$: $g(x)=(x-K)^+$ for the call,
$g\equiv1$ for the bond, $g(x)=\ind\{x>K\}$ for the digital \ts{13}{20}{14:46}.

\begin{finance}[Why call it $C_K(t,T)$]
Blythe is explicit that the notation separates the two roles of time: $K$ and $T$ are contract terms
fixed once and for all, while $t$ is the calendar date at which the price is observed
\ts{13}{20}{15:44}. Determining $C_K(t,T)$ for $t<T$ from a model is exactly the Black--Scholes problem
of \cref{ch:bs}; this chapter instead asks what can be said \emph{without} a model, from the family
$\{C_K(t,T)\}_{K}$ taken as given market data.
\end{finance}

The one assumption of the chapter is the Fundamental Theorem of Asset Pricing, already proved for a
finite state space in \cref{thm:ftap} (positive state prices $\Leftrightarrow$ no arbitrage, price
$=\bm\psi\T\bm C_T=\beta\,\E_\Q[\bm C_T]$). Its continuum version is:

\begin{theorem}[FTAP, continuum of states; not proved here]\label{ch22:thm-ftap}
Let $D(t,T)$ be the price at $t$ of a European claim with payout $g(S_T)$ at $T$. Under no arbitrage
(and suitable regularity, e.g.\ a complete diffusion market) there is a probability measure $\Q$, the
\emph{risk-neutral measure}, such that $D(t,T)/Z(t,T)$ is a $\Q$-martingale with respect to the
filtration generated by $S$; in particular, conditionally on $S_t$,
\begin{equation}\label{ch22:eq-ftap}
  D(t,T) = Z(t,T)\,\E_\Q\bigl[g(S_T)\mid S_t\bigr].
\end{equation}
\end{theorem}
\begin{proof}[Proof]
Not given here. \cref{thm:ftap} is the exact finite-dimensional statement and its complete proof
(Stiemke's lemma plus a separating hyperplane); the continuum, continuous-time version is due to
Harrison and Kreps (1979) and, as Blythe remarks, ``however many times you look at it, you're only
probably going to get through two or three pages before thinking, OK, that's hard''
\ts{13}{20}{51:57}. \cref{ch15:sec-onestep,ch15:sec-tree} prove it in the one-step and binomial-tree
cases, which already contain the essential idea: replicate the claim, then no-arbitrage forces its
price to equal the replicating portfolio's price, \cref{thm:ftap}(iii).
\end{proof}

\eqref{ch22:eq-ftap} is exactly \cref{thm:ftap}(iii), $C_0=\bm\psi\T\bm C_T=\beta\E_\Q[\bm C_T]$, with
the finite sum $\sum_i\psi_i(\cdot)$ replaced by the integral $\int(\cdot)\dd\Q$ and the state-price
vector $\bm\psi$ replaced by its continuum analogue: the density $\psi(x)=Z(t,T)f_\Q(x)$, where
$f_\Q$ is the $\Q$-density of $S_T$ given $S_t$, plays the role of the \emph{Arrow--Debreu price} of
state $x$ (\cref{sec:oneperiod} discusses the discrete Arrow--Debreu security $\bm e_i$ paying
\$1 in state $i$; here the ``state'' is a point $x\in\R_+$ and the security paying \$1 in a
neighbourhood $\dd x$ of $x$ costs $\psi(x)\dd x$). Applying \eqref{ch22:eq-ftap} to the digital
recovers the intuitive statement Blythe leads with: the price of \$1-if-the-coin-is-heads is the
probability of heads, discounted \ts{13}{20}{23:11}--\ts{13}{20}{24:48},
\begin{equation}\label{ch22:eq-digital-ftap}
  D_K(t,T) = Z(t,T)\,\Q(S_T>K\mid S_t).
\end{equation}

% ==================================================================
\section{From call prices to the density}\label{ch22:sec-bl-recap}

\subsection{The digital, twice}\label{ch22:sec-digital-twice}
The digital is priced two different ways: by \eqref{ch22:eq-digital-ftap} above, and as the limit of a
\emph{call spread}, exactly the computation of \cref{ch15:sec-digital}. Long $\lambda$ calls struck at
$K$ and short $\lambda$ calls struck at $K+1/\lambda$ has payout $0$ below $K$, slope $\lambda$ on
$[K,K+1/\lambda]$, and $1$ above $K+1/\lambda$ \ts{13}{20}{17:07}--\ts{13}{20}{19:19}; letting
$\lambda\to\infty$ this payout tends pointwise to $\ind\{S_T>K\}$, and its price tends to
\begin{equation}\label{ch22:eq-cs-limit}
  \lim_{\lambda\to\infty}\lambda\bigl[C_K(t,T)-C_{K+1/\lambda}(t,T)\bigr] = -\frac{\partial C_K}{\partial K}(t,T).
\end{equation}
This is \eqref{ch15:eq-digital-limit} with $\varepsilon=1/(2\lambda)$ and a one-sided spread instead of
a centred one; the limit is the same. Equating \eqref{ch22:eq-cs-limit} with
\eqref{ch22:eq-digital-ftap} gives
\begin{equation}\label{ch22:eq-cdf}
  \Q(S_T>K\mid S_t) = -\frac{1}{Z(t,T)}\frac{\partial C_K}{\partial K}(t,T),
  \qquad
  F_\Q(K):=\Q(S_T\le K\mid S_t) = 1+\frac{1}{Z(t,T)}\frac{\partial C_K}{\partial K}(t,T),
\end{equation}
and one more strike derivative gives the density, restating \cref{ch03:cor-bl} with the discount factor
made explicit:
\begin{equation}\label{ch22:eq-bl}
  f_\Q(K) = \frac{1}{Z(t,T)}\frac{\partial^2 C_K}{\partial K^2}(t,T).
\end{equation}
\eqref{ch22:eq-bl} is the \emph{Breeden--Litzenberger relation}, already proved in
\cref{ch03:cor-bl} (there for $Z\equiv1$, i.e.\ undiscounted expectations) and used in
\cref{ch15:sec-digital} for the digital alone; nothing here is new mathematics, only the
observation, which is the point of the lecture, that \emph{iterating the same one-strike-derivative
trick that gives the digital, one more time, gives the entire density, not just one tail probability}.

\begin{intuition}[One relation, no model]
Nothing about $\Q$ or $S$ entered the derivation beyond FTAP itself: \eqref{ch22:eq-bl} holds for
\emph{whatever process} generates the risk-neutral distribution, log-normal (Black--Scholes) or not
\ts{13}{20}{39:03}--\ts{13}{20}{39:31}. Plugging the Black--Scholes formula \eqref{ch15:eq-call} into
$\partial^2_KC$ recovers the log-normal density of \cref{ch15:eq-gbm}, but \eqref{ch22:eq-bl} is
agnostic to that choice: it converts \emph{observed} call prices into a distribution without ever
positing a model for $S$.
\end{intuition}

\subsection{The call butterfly: trading the density}\label{ch22:sec-butterfly}
A single strike derivative (the call spread) has an obvious trading analogue; Blythe's second
construction shows that the \emph{second} strike derivative does too \ts{13}{20}{30:56}. Consider the
portfolio, long $\lambda$ calls at $K-1/\lambda$, short $2\lambda$ calls at $K$, long $\lambda$ calls at
$K+1/\lambda$ (a \emph{call butterfly}, traded in practice, unlike the idealised limit of
\cref{ch22:sec-digital-twice} \ts{13}{20}{32:22}--\ts{13}{20}{32:37}). Its price is
\begin{equation}\label{ch22:eq-butterfly-price}
  B_K(\lambda;t,T) := \lambda\Bigl[C_{K-1/\lambda}(t,T) - 2C_K(t,T) + C_{K+1/\lambda}(t,T)\Bigr],
\end{equation}
a centred second difference of the call price, scaled by $\lambda$ (one more power of $\lambda$ than
the call spread, matching the second derivative being approximated by
$C(K-\varepsilon)-2C(K)+C(K+\varepsilon)\approx\varepsilon^2C''(K)$ with $\varepsilon=1/\lambda$):
\begin{equation}\label{ch22:eq-butterfly-limit}
  \lim_{\lambda\to\infty}\lambda\,B_K(\lambda;t,T) = \frac{\partial^2 C_K}{\partial K^2}(t,T) = Z(t,T)f_\Q(K).
\end{equation}

\begin{proposition}[The butterfly payoff is a triangle]\label{ch22:prop-triangle}
The payout of the portfolio in \eqref{ch22:eq-butterfly-price} at maturity is the triangular tent
function of base $[K-1/\lambda,K+1/\lambda]$ and height $1$ at $K$,
\[
  h_\lambda(x) = \lambda\Bigl[(x-K+\tfrac1\lambda)^+ - 2(x-K)^+ + (x-K-\tfrac1\lambda)^+\Bigr]
  = \bigl(1-\lambda|x-K|\bigr)^+ .
\]
\end{proposition}
\begin{proof}
Direct case check on the three pieces of the piecewise-linear function; each summand is a hockey stick
with a kink at $K\mp1/\lambda$ or $K$, and the three slopes ($\lambda$, $-2\lambda$, $\lambda$) sum to
$0$ outside $[K-1/\lambda,K+1/\lambda]$ (so the payout is $0$ there) and change by $\mp\lambda$ at each
kink, giving slope $+\lambda$ on $[K-1/\lambda,K]$ and $-\lambda$ on $[K,K+1/\lambda]$, height
$1=\lambda\cdot\frac1\lambda$ at $K$. This is exactly the sketch in class
\ts{13}{20}{36:17}--\ts{13}{20}{36:41}, made explicit.
\end{proof}

This makes \eqref{ch22:eq-butterfly-limit} the exact analogue of the elementary fact
$\Q(K-\delta<S_T<K+\delta)\approx2\delta f_\Q(K)$ for small $\delta$: the triangle of
\cref{ch22:prop-triangle} has base $2/\lambda$ and height $1$, hence area $1/\lambda$, and by FTAP its
discounted expected payout is
\begin{equation}\label{ch22:eq-triangle-area}
  B_K(\lambda;t,T) = Z(t,T)\,\E_\Q[h_\lambda(S_T)\mid S_t] \approx Z(t,T)\cdot\frac1\lambda\, f_\Q(K)
  \quad\text{for large }\lambda,
\end{equation}
which is \eqref{ch22:eq-butterfly-limit} divided by $\lambda$ on both sides
\ts{13}{20}{35:39}--\ts{13}{20}{38:16}. Unlike the call spread, the butterfly is genuinely traded
(interest-rate desks quote butterflies $10$--$20$ basis points wide, as noted in class
\ts{13}{20}{33:24}--\ts{13}{20}{33:37}), so \eqref{ch22:eq-triangle-area} is not a theoretical limit
alone: a market maker who believes $\Q(S_T\approx K)$ is larger than the price of the corresponding
butterfly implies will buy it \ts{13}{20}{38:23}--\ts{13}{20}{38:48}.

\begin{table}[ht]
\centering
\caption{Butterfly price $\lambda\,B_K(\lambda;0,T)$ converging to $C''(K)=e^{-rT}f_\Q(K)$, at $K=80$,
Black--Scholes with $S_0=80$, $r=5\%$, $\sigma=20\%$, $T=2$ (parameters of \cref{ch15:ex-atm}).
Computed for these notes.}\label{ch22:tab-butterfly}
\begin{tabular}{@{}rrr@{}}
\toprule
$\lambda$ & $\varepsilon=1/\lambda$ & $\lambda\,B_{80}(\lambda)$\\
\midrule
$0.05$ & $20$   & $0.014656$\\
$0.1$  & $10$   & $0.015354$\\
$0.2$  & $5$    & $0.015536$\\
$0.5$  & $2$    & $0.015588$\\
$1$    & $1$    & $0.015596$\\
$2$    & $0.5$  & $0.015598$\\
$5$    & $0.2$  & $0.015598$\\
\midrule
\multicolumn{2}{r}{exact $e^{-rT}f_\Q(80)$} & $0.015598$\\
\bottomrule
\end{tabular}
\end{table}

\begin{figure}[ht]
\centering
\begin{tikzpicture}
\begin{axis}[width=0.92\linewidth,height=6cm,axis lines=left,xmin=35,xmax=145,ymin=0,ymax=0.017,
  xlabel={$K$},ylabel={discounted density},tick label style={font=\footnotesize},
  label style={font=\small},
  legend style={at={(0.97,0.97)},anchor=north east,font=\scriptsize,draw=none,fill=none}]
  \addplot[thick,black!45,mark=none] coordinates {
   (40.0,0.00092) (44.0,0.00194) (48.0,0.00347) (52.0,0.00545) (56.0,0.00770) (60.0,0.00999)
   (64.0,0.01208) (68.0,0.01377) (72.0,0.01494) (76.0,0.01554) (80.0,0.01560) (84.0,0.01518)
   (88.0,0.01439) (92.0,0.01333) (96.0,0.01211) (100.0,0.01081) (104.0,0.00950) (108.0,0.00824)
   (112.0,0.00707) (116.0,0.00600) (120.0,0.00504) (124.0,0.00421) (128.0,0.00349) (132.0,0.00287)
   (136.0,0.00235) (140.0,0.00192)};
  \addplot[only marks,mark=o,mark size=1.4pt,defcol] coordinates {
   (40.0,0.00092) (44.0,0.00194) (48.0,0.00347) (52.0,0.00545) (56.0,0.00770) (60.0,0.00999)
   (64.0,0.01208) (68.0,0.01377) (72.0,0.01494) (76.0,0.01554) (80.0,0.01560) (84.0,0.01518)
   (88.0,0.01439) (92.0,0.01333) (96.0,0.01211) (100.0,0.01081) (104.0,0.00950) (108.0,0.00824)
   (112.0,0.00707) (116.0,0.00600) (120.0,0.00504) (124.0,0.00421) (128.0,0.00349) (132.0,0.00287)
   (136.0,0.00235) (140.0,0.00192)};
  \legend{$e^{-rT}f_\Q$ (exact),{$\lambda B_K(\lambda)$, $\lambda=8$}}
\end{axis}
\end{tikzpicture}
\caption{Discounted risk-neutral density recovered from the call-butterfly price
\eqref{ch22:eq-butterfly-price} at $\lambda=8$ (width $2/\lambda=0.25$), against the exact
Black--Scholes log-normal density (same parameters as \cref{ch22:tab-butterfly}); markers overlap the
line to plotting accuracy. Computed for these notes.}\label{ch22:fig-density}
\end{figure}

% ==================================================================
\section{From the density to any price: FTAP the other way}\label{ch22:sec-density-to-price}

\cref{ch22:sec-bl-recap} went from call prices to a distribution. The Fundamental Theorem
\eqref{ch22:eq-ftap} already goes the other way, and the resulting statement is worth writing out
explicitly \ts{13}{20}{49:41}--\ts{13}{20}{54:09}: if $f_\Q(\cdot\mid S_t)$ is known, the price of
\emph{any} European claim with payout $g$ is
\begin{equation}\label{ch22:eq-lotus}
  D(t,T) = Z(t,T)\int_0^\infty g(x)\,f_\Q(x\mid S_t)\,\dd x,
\end{equation}
an application of the law of the unconscious statistician (LOTUS): $\E_\Q[g(S_T)]$ is computed by
integrating $g$ against the density of $S_T$, without first working out the (generally intractable)
density of the transformed variable $g(S_T)$ itself. Taking $g(x)=(x-K)^+$ in \eqref{ch22:eq-lotus}
recovers \cref{ch03:thm-call} (there stated for a general random variable $X$ under the physical
measure; here $X=S_T$ under $\Q$, with the discount factor $Z(t,T)$ attached). So the two directions
of the duality are: call prices $\to$ density via two strike derivatives
(\cref{ch22:sec-bl-recap}), density $\to$ any European price via one integral
(\eqref{ch22:eq-lotus}), and both use nothing beyond FTAP.

% ==================================================================
\section{Calls span every European payoff}\label{ch22:sec-spanning}

The two directions combine into a genuinely new statement, not reducible to
Breeden--Litzenberger alone: knowing $C_K(t,T)$ for every $K$ determines the price of \emph{every}
European claim, and moreover gives an \emph{explicit static replicating portfolio} built only from
the zero-coupon bond, the stock, and calls \ts{13}{20}{55:44}--\ts{13}{20}{56:09}.

\subsection{Piecewise-linear payouts}\label{ch22:sec-pwl}
If $g$ is piecewise linear with kinks at $K_1<K_2<\dots<K_n$, its graph is matched exactly, for
$x\ge K_1$ say, by a portfolio holding $\theta_i$ calls of strike $K_i$, where $\theta_i$ is the
change in slope of $g$ at $K_i$ (a call of strike $K_i$ contributes slope $1$ for $x>K_i$ and $0$
below, so a linear combination reproduces any sequence of slope changes)
\ts{13}{20}{56:47}--\ts{13}{20}{58:13}. Since $g$ and this call portfolio have \emph{identical} payoff
at $T$ for every value of $S_T$, no-arbitrage (\cref{thm:ftap}(iii): a replicable claim has a unique
price, the price of any replicating portfolio) forces their prices to agree at every $t\le T$,
\begin{equation}\label{ch22:eq-pwl-repl}
  D(t,T) = \sum_i \theta_i\, C_{K_i}(t,T), \qquad \theta_i=\Delta g'(K_i).
\end{equation}
A continuous $g$ is a uniform limit of piecewise-linear functions, so \eqref{ch22:eq-pwl-repl} suggests
(but does not by itself prove, because of the limit) that \emph{every} continuous payout is spanned by
calls. \cref{ch22:sec-taylor} makes this precise and quantitative for $g\in C^2$, which is the
practically relevant case and the one the lecture treats.

\begin{coursenote}[A digression worth keeping]
A student asks whether replication at $T$ really forces equality of prices at $t<T$ for maturities
decades away, given that no one can actually hold the position without funding or margin risk in the
meantime \ts{13}{20}{59:44}--\ts{13}{20}{1:02:13}. Blythe's answer: mathematically yes, under
no-arbitrage; empirically, for very long maturities the two prices \emph{have} been observed to drift
apart, because ``no arbitrage'' assumes unlimited risk capital willing to hold the position for the
full horizon, which breaks down since 2008. This is a caveat on the economic content of
\cref{thm:ftap}(iii), not on its mathematics, and it is not pursued further here.
\end{coursenote}

\subsection{The exact Carr--Madan identity}\label{ch22:sec-taylor}
\begin{proposition}[Static replication of a smooth payoff]\label{ch22:prop-carrmadan}
Let $g:\R_+\to\R$ be twice continuously differentiable. For every $x\ge0$,
\begin{equation}\label{ch22:eq-taylor-exact}
  g(x) = g(0) + g'(0)\,x + \int_0^\infty g''(K)\,(x-K)^+\,\dd K.
\end{equation}
\end{proposition}
\begin{proof}
This is the lecture's ``exact Taylor series to second order'' \ts{13}{20}{1:03:33}--\ts{13}{20}{1:04:20},
here proved rather than only asserted. Start from the fundamental theorem of calculus twice,
\[
  g(x) = g(0) + \int_0^x g'(K)\,\dd K, \qquad g'(K) = g'(0) + \int_0^K g''(c)\,\dd c,
\]
substitute and swap the order of integration in the double integral
$\int_0^x\!\int_0^K g''(c)\,\dd c\,\dd K$ over the triangle $0\le c\le K\le x$:
\[
  \begin{gathered}
  \int_0^x\int_0^K g''(c)\,\dd c\,\dd K = \int_0^x g''(c)\int_c^x \dd K\,\dd c = \int_0^x g''(c)(x-c)\,\dd c \\
  = \int_0^\infty g''(c)\,(x-c)^+\,\dd c,
  \end{gathered}
\]
where the last equality extends the integral to $(x,\infty)$ for free because $(x-c)^+=0$ there. Renaming
the dummy variable $c\to K$ gives \eqref{ch22:eq-taylor-exact}. (Equivalently, integrate the right-hand
side of \eqref{ch22:eq-taylor-exact} by parts twice, as the lecture suggests
\ts{13}{20}{1:04:20}--\ts{13}{20}{1:04:31}, which reverses this computation.)
\end{proof}

\begin{remark}
The identity holds pointwise, for \emph{every} $x\ge0$, with no approximation and no limit: this is the
sense in which it is ``exact''. It is the content of \cref{ch22:sec-pwl} made rigorous:
$g''(K)\dd K$ replaces the discrete slope changes $\theta_i=\Delta g'(K_i)$ of a piecewise-linear
function by a continuum of infinitesimal slope changes, each priced by $(x-K)^+$, the call payoff.
Where the lecture stops after the piecewise-linear case and asserts the smooth case by continuity
\ts{13}{20}{1:02:52}--\ts{13}{20}{1:03:03}, \cref{ch22:prop-carrmadan} needs no limiting argument at
all: it is an identity between functions, true for every fixed $x$, and the piecewise-linear intuition
of \cref{ch22:sec-pwl} is recovered by taking $g''$ to be a sum of Dirac masses (formally, outside the
$C^2$ hypothesis, but the calls-with-strikes-at-the-kinks picture is exactly this limit).
\end{remark}

Set $x=S_T$ in \eqref{ch22:eq-taylor-exact} (a pointwise identity between random variables, valid on
every sample path) and apply the discounted-expectation operator $Z(t,T)\,\E_\Q[\cdot\mid S_t]$ to both
sides. The left side becomes $D(t,T)$, the price of the claim with payout $g$. On the right,
linearity of expectation gives three terms: $g(0)$ times the price of $g(0)$ zero-coupon bonds is
$Z(t,T)g(0)$; the term $g'(0)\,\E_\Q[S_T\mid S_t]$, discounted, is $g'(0)$ times the price of \emph{one
unit of stock} (the stock, viewed as its own European claim with payout $S_T$ and replicating portfolio
``hold the stock'', so $Z(t,T)\E_\Q[S_T\mid S_t]=S_t$ by \cref{ch22:thm-ftap} applied to $g(x)=x$
\ts{13}{20}{1:07:25}--\ts{13}{20}{1:07:45}); and the integral term becomes, exchanging expectation and
the $K$-integral (justified when $g''\ge0$ by Tonelli, or under an integrability bound on $g''$ in
general),
\[
  \begin{gathered}
  Z(t,T)\,\E_\Q\Bigl[\int_0^\infty g''(K)(S_T-K)^+\dd K\ \Big|\ S_t\Bigr]
  = \int_0^\infty g''(K)\,Z(t,T)\,\E_\Q\bigl[(S_T-K)^+\mid S_t\bigr]\dd K \\
  = \int_0^\infty g''(K)\,C_K(t,T)\,\dd K.
  \end{gathered}
\]
Collecting the three terms:

\begin{theorem}[Calls span all European payoffs; Carr--Madan]\label{ch22:thm-spanning}
For $g\in C^2(\R_+)$ with the payoff at maturity $g(S_T)$,
\begin{equation}\label{ch22:eq-spanning}
  D(t,T) = g(0)\,Z(t,T) + g'(0)\,S_t + \int_0^\infty g''(K)\,C_K(t,T)\,\dd K.
\end{equation}
Equivalently, the claim is statically replicated at $t$ by $g(0)$ zero-coupon bonds, $g'(0)$ units of
stock, and $g''(K)\dd K$ calls of strike $K$ for every $K>0$.
\end{theorem}

\begin{keyformulas}
\textbf{Call prices $\to$ distribution:} $\displaystyle \Q(S_T>K\mid S_t)=-\frac1{Z}\partial_KC_K$,
\quad $\displaystyle f_\Q(K)=\frac1{Z}\partial^2_KC_K$ \eqref{ch22:eq-cdf}--\eqref{ch22:eq-bl}
(Breeden--Litzenberger, \cref{ch03:cor-bl}). \\
\textbf{Distribution $\to$ any price:} $\displaystyle D(t,T)=Z(t,T)\int g(x)f_\Q(x\mid S_t)\dd x$
\eqref{ch22:eq-lotus} (FTAP/LOTUS). \\
\textbf{Static replication (Carr--Madan):}
$\displaystyle D(t,T)=g(0)Z(t,T)+g'(0)S_t+\int_0^\infty g''(K)C_K(t,T)\dd K$ \eqref{ch22:eq-spanning}.\\
\textbf{Traded proxy for the density:} call butterfly $\lambda B_K(\lambda)\to Z(t,T)f_\Q(K)$
\eqref{ch22:eq-butterfly-limit}.
\end{keyformulas}

\begin{example}[Checking the identity on $g(x)=x^2$]\label{ch22:ex-x2}
Take $g(x)=x^2$, so $g(0)=0$, $g'(0)=0$, $g''(K)=2$ for all $K$. \eqref{ch22:eq-taylor-exact} claims
$x^2=\int_0^\infty 2(x-K)^+\dd K$, and indeed $\int_0^x 2(x-K)\dd K=2\bigl[xK-K^2/2\bigr]_0^x=x^2$: a
direct check of \cref{ch22:prop-carrmadan} with an elementary integral. Numerically (trapezoidal rule,
grid $[0,800]$, $4\times10^5$ points, computed for these notes) the same integral gives $6400.0000$
against the exact $80^2=6400$. A second check with a bounded, non-polynomial $g(x)=e^{-0.01x}$
($g''(K)=10^{-4}e^{-0.01K}$) gives, at $x=80$, exact $g(80)=0.449329$ against the numerical
reconstruction $0.449329$: \eqref{ch22:eq-taylor-exact} is confirmed to $6$ digits both for a payoff
whose second derivative is constant and for one that decays.
\end{example}

% ==================================================================
\section{What the terminal distribution does not determine}\label{ch22:sec-process}

Everything above concerns a single fixed maturity $T$: knowing all call prices $C_K(t,T)$ for that one
$T$ determines the distribution of $S_T$ given $S_t$, and hence the price of every payoff that depends
on $S_T$ alone. The natural next question is whether knowing this for \emph{every} maturity $T$ (a full
implied-volatility surface, in market language) determines the entire stochastic process $(S_u)_{u\ge
t}$ \ts{13}{20}{1:12:11}--\ts{13}{20}{1:13:07}.

\begin{proposition}[Marginals do not determine the process; not proved here]\label{ch22:prop-dupire}
Knowledge of the marginal (risk-neutral) distribution of $S_T$ given $S_t$, for every $T\ge t$, does
not determine the law of the process $(S_u)_{u\in[t,T]}$. If in addition $S$ is assumed to be a
one-dimensional diffusion driven by Brownian motion, $\dd S_u=\mu(u,S_u)\dd u+\sigma(u,S_u)\dd W_u$,
then the marginals \emph{do} determine the process, via the local volatility function
$\sigma_{\mathrm{loc}}(u,x)$ (Dupire, Derman--Kani, early 1990s).
\end{proposition}
\begin{proof}[Discussion, not a proof]
Blythe gives only the intuition, not the construction, and none is attempted here (it belongs to
stochastic calculus proper, beyond \cref{ch:sde}): take a denser and denser grid of maturities with
prescribed, say symmetric, marginals that spread out smoothly in $T$; the stock can still be made to
jump discontinuously between two adjacent grid times while landing in a distribution consistent with
the (denser and denser) marginals at both ends, because \emph{nothing in the one-maturity-at-a-time
data constrains the joint law across maturities} \ts{13}{20}{1:13:34}--\ts{13}{20}{1:14:07}. Ruling
out such jumps, by assuming a diffusion, is what lets Dupire's formula reconstruct a unique
local-volatility function from the full surface of \eqref{ch22:eq-bl}-derived marginals.
\end{proof}

\begin{finance}[Why this matters for a trading desk]
A desk that has priced, and can perfectly hedge with the observed vanilla surface, every European
payoff via \eqref{ch22:eq-spanning}, still cannot claim the same for a path-dependent exotic (a
barrier, an Asian option, or, as Blythe's example, a derivative paying on the number of times $S$
crosses a level in a small interval \ts{13}{20}{1:19:05}--\ts{13}{20}{1:19:10}): its price is not
pinned down by the vanilla surface alone, precisely because \cref{ch22:prop-dupire} says the surface
does not pin down the process. A student's question makes the link to market completeness explicit
\ts{13}{20}{1:16:44}--\ts{13}{20}{1:18:52}: a jump/flip process for which the marginals do not
determine the law corresponds to an \emph{incomplete} market (\cref{def:onepmodel}) with respect to
the vanilla calls, exactly as more states than independent assets leave state prices non-unique in the
one-period model of \cref{sec:oneperiod}. The risk-neutral density $f_\Q(\cdot\mid S_t)$ at each
fixed $T$ is still uniquely pinned down by the calls (Blythe is emphatic on this point when pressed
\ts{13}{20}{1:19:17}--\ts{13}{20}{1:19:48}); what is not determined is the joint law across different
$T$'s, i.e.\ the path.
\end{finance}

\begin{remark}[An alternative spanning set: Arrow--Debreu securities]
A question from the audience asks whether calls are the \emph{only} convenient spanning set
\ts{13}{20}{1:16:16}--\ts{13}{20}{1:17:41}. They are not: knowing all digitals gives the CDF directly
(one derivative instead of two), and knowing the price of every Arrow--Debreu security $\ind\{S_T\in
\dd x\}$ gives the density directly, with no derivative at all. Blythe's preference for calls is that
the identity \eqref{ch22:eq-spanning} is unusually clean, expressing an arbitrary smooth payoff through
an explicit, liquidly-traded static hedge, whereas Arrow--Debreu securities of infinitesimal width are
not traded and digitals are less liquid than calls (\cref{ch15:sec-digital}). This is also exactly the
representation \cref{ch:ross} will use (state prices as a discrete matrix, recovered via
Perron--Frobenius, \cref{sec:pf}) in a setting where the underlying state space is finite rather than
a continuum of strikes.
\end{remark}

% ==================================================================
\section*{Exercises}
\addcontentsline{toc}{section}{Exercises}
All exercises below are (not from the course).

\begin{exercise}[Butterfly width and bias]\label{ch22:ex-width}
For the Black--Scholes parameters of \cref{ch22:tab-butterfly} ($S_0=80$, $r=5\%$, $\sigma=20\%$,
$T=2$), let $\lambda=4$ (butterfly of half-width $1/\lambda=0.25$ around each wing, full base
$0.5$). Compute $B_{80}(4)$ from the closed-form Black--Scholes call price \eqref{ch15:eq-call} and
compare $\lambda B_{80}(4)$ to the exact value in \cref{ch22:tab-butterfly}. Express the leading-order
error as a function of $\lambda$ using a fourth-order Taylor expansion of $C$ around $K=80$, and check
it is consistent with the pattern in the table.
\end{exercise}

\begin{exercise}[Straddle via Carr--Madan]\label{ch22:ex-straddle}
A straddle pays $g(S_T)=|S_T-K_0|$ for a fixed $K_0>0$. This $g$ is not $C^2$ at $K_0$ (it has a
kink), so \eqref{ch22:eq-taylor-exact} must be applied with $g''$ interpreted as a Dirac mass
$2\delta_{K_0}$ (the jump in $g'$ at $K_0$ is $2$). Using \eqref{ch22:eq-pwl-repl} directly instead
(the straddle is piecewise linear with one kink), find the static replicating portfolio of zero-coupon
bonds, stock and calls of a single strike, and check it agrees with formally substituting $g''=2\delta_{K_0}$
into \eqref{ch22:eq-spanning}.
\end{exercise}

\begin{exercise}[Digital from a butterfly, not a call spread]\label{ch22:ex-digital-butterfly}
Show that a digital option cannot be replicated as a \emph{limit of call butterflies} the way it was
replicated as a limit of call spreads in \cref{ch22:sec-digital-twice}: compute
$\lim_{\lambda\to\infty}B_K(\lambda;t,T)$ (not $\lambda B_K(\lambda;t,T)$) and interpret the result in
terms of the payoff triangle of \cref{ch22:prop-triangle} shrinking to a point.
\end{exercise}

\begin{exercise}[Two maturities, same terminal law, different paths]\label{ch22:ex-two-paths}
Let $W$ be a Brownian motion and $Z\sim N(0,1)$. Show that $S_t=1+W_t$ and $\tilde S_t=1+\sqrt t\,Z$
have the same marginal law $N(1,t)$ at every $t\in[0,2]$, but different laws as processes: compute
$\Cov(S_1,S_2)$ and $\Cov(\tilde S_1,\tilde S_2)$. Which of the two is a martingale? This illustrates
\cref{ch22:prop-dupire} concretely: a claim paying $(S_2-S_1)^2$ has different prices under the two
processes although all the terminal laws, hence all the call prices, agree.
\end{exercise}
